(1) 谐振子基态波函数为 \(\psi_0(x)=(\pi a^2)^{-1/4}e^{-x^2/(2a^2)}\)，归一化由高斯积分保证。故原点密度
\[
|\psi_0(0)|^2=\frac1{\sqrt\pi\,a}.
\]

(2) 定态非简并微扰论的一阶能量校正为 \(\Delta E_n^{(1)}=\langle n|V|n\rangle=\lambda|\psi_n(0)|^2\)。因此
\[
\Delta E_0^{(1)}=\frac{\lambda}{\sqrt\pi\,a}.
\]
\(\lambda >0\) 的排斥 \(\delta \) 势提高能量，\(\lambda <0\) 的吸引势降低能量。

(3) 一维谐振子波函数宇称为 \((-1)^n\)，所有奇数 n 态 \(\psi_n(0)=0\)，故
\[
\Delta E_{2j+1}^{(1)}=0.
\]
此结论只说明一阶能移为零；更高阶微扰仍须具体核算，但原点处奇态本身因宇称始终有节点。

(4) n=2 波函数 \(\psi_2(x)=\psi_0(x)(2(x/a)^2-1)/\sqrt2\)（按 Hermite 多项式归一），故 \(\psi_2(0)=-\psi_0(0)/\sqrt2\)。因此
\[
|\psi_2(0)|^2=\frac1{2\sqrt\pi\,a},\quad
\Delta E_2^{(1)}=\frac{\lambda}{2\sqrt\pi\,a}
=\frac12\Delta E_0^{(1)}.
\]
